\documentclass[11pt]{article}
\usepackage[margin=1in]{geometry}
\usepackage{amsmath,amssymb,amsthm}
\usepackage{booktabs}
\usepackage[colorlinks=true,linkcolor=blue,citecolor=blue,urlcolor=blue]{hyperref}

\newtheorem{theorem}{Theorem}[section]
\newtheorem{lemma}[theorem]{Lemma}
\newtheorem{proposition}[theorem]{Proposition}
\newtheorem{conjecture}[theorem]{Conjecture}
\newtheorem{question}[theorem]{Question}
\theoremstyle{definition}
\newtheorem{finding}[theorem]{Finding}
\theoremstyle{remark}
\newtheorem{remark}[theorem]{Remark}

\newcommand{\Z}{\mathbb{Z}}
\newcommand{\F}{\mathbb{F}}
\newcommand{\cP}{\mathcal{P}}
\DeclareMathOperator{\lcm}{lcm}
\DeclareMathOperator{\ord}{ord}
\setlength{\emergencystretch}{1.5em}

\title{On the Infeasibility of Constructing Lenstra--Pomerance
Counterexamples to Agrawal's Conjecture\\
via Twin-Smooth Subset Products}
\author{Dan Shumow \and Claude (Anthropic Fable~5)}
\date{August 24, 2026}

\begin{document}
\maketitle

\begin{abstract}
Agrawal's conjecture asserts that for $r$ prime, $r \nmid n$, the congruence
$(X-1)^n \equiv X^n-1 \pmod{n,\,X^r-1}$ forces $n$ prime or $n^2 \equiv 1
\pmod r$; if true it yields a deterministic $\tilde O(\log^3 n)$ primality
test. The only known route to a counterexample is the Lenstra--Pomerance
construction: a squarefree $n = p_1\cdots p_k$ that is simultaneously
Carmichael and Lucas--Carmichael, with side conditions, at a size heuristics
place at hundreds of digits or (far) beyond --- out of reach of search, but
potentially in reach of \emph{construction}. We report on a systematic attempt
to execute that construction as a two-stage pipeline: harvest a pool of
admissible primes (twin-smooth integers with prime sum in a fixed congruence
class), then solve a two-sided 0/1 subset-product problem over
$(\Z/\Lambda_-)^\times \times (\Z/\Lambda_+)^\times$. The outcome is
asymmetric. The \emph{supply} stage is solved: a partition-first harvest
sidesteps a measured combinatorial obstruction (a universal
$\approx 0.06$ survivors-per-demand-bit invariant of partition-agnostic twin
corpora, measured up to the complete $B=547$ corpus of $82{,}026{,}426$
pairs) and delivers validated pool-to-demand ratios of $13.8$--$44$, an
over-provisioning of one to two orders of magnitude. The \emph{solver} stage
is blocked: generic exact and heuristic solvers cap two to three orders of
magnitude below the required odd-part dimension, and the one known scalable
family --- the subexponential subset-product methods of
Alford--Grantham--Hayman--Shallue (AGHS) and their generalized-birthday core
--- requires pool primes with high \emph{identity density} on the group
components. We quantify both jaws of the resulting vise. Under the double
Carmichael + Lucas--Carmichael condition, the divisor-paradigm harvest that
would supply AGHS-like density yields an essentially empty pool ($0$--$2$
primes) at all feasible scales; and a direct measurement of birthday-solver
reach gives the law $r_{\max} \approx O(\log N)/(1-\rho)$, so spanning our
thousand-component groups needs identity density $\rho \gtrsim 0.98$, while
harvestable pools deliver $\rho \approx 0.003$--$0.008$ and enlarging the
harvest improves reach only logarithmically. High-density pools are
unharvestable; low-density pools are unsolvable. This two-part quantified
obstruction closes the $r=5$ Lenstra--Pomerance route for the entire
AGHS/birthday solver class. We emphasize scope: this is not a disproof of
Agrawal's conjecture, nor a proof that no counterexample can be constructed;
the directions left open are a fundamentally different solver paradigm and a
different AKS modulus $r$ (a different construction).
\end{abstract}

\section{Introduction}\label{sec:intro}

For coprime integers $n, r$ write $T(-1,n,r)$ for the congruence
\begin{equation}\label{eq:aks}
(X-1)^n \equiv X^n - 1 \pmod{n,\, X^r-1}.
\end{equation}

\begin{conjecture}[Agrawal; Bhattacharjee--Pandey \cite{BP01}, popularized
in the AKS paper \cite{AKS}]\label{conj:agrawal}
If $r$ is a prime not dividing $n$ and $T(-1,n,r)$ holds, then $n$ is prime or
$n^2 \equiv 1 \pmod r$.
\end{conjecture}

If Conjecture~\ref{conj:agrawal} were true, a single evaluation of
\eqref{eq:aks} at one $r = O(\log n)$ would certify primality, giving a
deterministic $\tilde O(\log^3 n)$ test, versus $\tilde O(\log^6 n)$ for the
best proven algorithm \cite{LP2019}. Its status is strange: Lenstra and
Pomerance \cite{LPremarks} gave a heuristic argument for
$\gg x^{1-\varepsilon}$ counterexamples below $x$, yet the conjecture has
survived exhaustive search to $10^{17}$ (Primaboinca, 2010--2020) and
folklore places the least counterexample at hundreds of digits. The only
concrete counterexample route --- the Lenstra--Pomerance (L--P) construction
(Theorem~\ref{thm:lp} below), extended by Popovych \cite{Popovych09} and
Hegde--Devaraj \cite{HD21} --- demands an $n$ that is simultaneously a
Carmichael and a Lucas--Carmichael number with additional congruence
conditions; no such number is known, and none exists below $10^{18}$
\cite{Pinch}.

\emph{The constructive program.} Since any counterexample refutes the
conjecture, minimality can be abandoned. The L--P density heuristic
\emph{improves} with size, and the construction is deterministic: when its
arithmetic side conditions hold, \eqref{eq:aks} follows with no further luck
required. This suggests a pipeline in the style of the large-scale Carmichael
constructions of Alford--Granville--Pomerance \cite{AGP} and L\"oh--Niebuhr
\cite{LN}: \textbf{harvest} a pool $\cP$ of primes satisfying the local
conditions --- which turn out to be exactly the \emph{twin smooth integers
with prime sum} of isogeny-based cryptography
\cite{BSIDH,CMN21,BCCEMNS23,BHLNV,MSW} in a fixed congruence class --- and
then \textbf{solve} for a subset $S \subseteq \cP$ whose product satisfies a
two-sided congruence, a 0/1 subset-product problem in the group
$(\Z/\Lambda_-)^\times \times (\Z/\Lambda_+)^\times$
(Proposition~\ref{prop:reduction}). This paper reports what happened when
that pipeline was built and measured end to end.

\emph{Contributions.}
\begin{enumerate}\itemsep2pt
\item[(i)] \textbf{A validated, over-provisioned supply method}
(\S\ref{sec:supply}). We first measure a universal obstruction to mining
existing twin-smooth corpora: across bounds $B = 100$--$547$ and pool sizes
$78$ to $73{,}163$ (the complete $B=547$ corpus of $82{,}026{,}426$ pairs),
the best side-partition retains only $\approx 0.06$ pool elements per bit of
group demand --- a factor $16$--$25$ below feasibility. We then show that a
\emph{partition-first} harvest sidesteps this invariant entirely, with
measured pool-to-demand ratios of $13.8$--$44$ under a yield model validated
(conservatively) against the complete $B=547$ corpus.
\item[(ii)] \textbf{Identification and quantification of a two-part
solver-stage obstruction} (\S\ref{sec:solver}). The subset-product problem
splits into an easy $2$-part ($\approx 16\%$ of demand bits, solved exactly
at full scale) and a dominant odd part ($\approx 84\%$) on which generic
exact and heuristic solvers cap two to three orders of magnitude short. The
scalable family --- the AGHS $\omega$-guided reduction \cite{AGHS}, which
built a Carmichael number with $10{,}333{,}229{,}505$ prime factors, and the
generalized-birthday machinery \cite{Wagner} at its core --- needs pools of
high \emph{identity density}. We quantify both jaws of the resulting vise.
\emph{Harvest jaw:} the divisor-paradigm harvest that would achieve
AGHS-like density ($\approx 0.3$--$0.5$) is empty ($0$--$2$ primes) at every
feasible scale, while harvestable pools have density
$\approx 0.003$--$0.008$. \emph{Solver jaw:} a correct Wagner-$4$ birthday
solver, measured across density and pool size, obeys the reach law
$(1-\rho)\, r_{\max} \approx \mathrm{const}$ with $r_{\max}$ growing only
logarithmically in the pool --- so spanning our thousand-component groups
requires $\rho \gtrsim 0.98$, two-plus orders of magnitude above what any
harvestable pool supplies, and harvesting more cannot close the gap. The
pipeline is blocked at the solver: high-density pools are unharvestable,
low-density pools are unsolvable.
\item[(iii)] \textbf{Reusable machinery and diagnostics}
(\S\ref{sec:conclusion}): a validated twin-smooth yield model, a
near-complete C implementation of the CHM closure, a solver-frontier map for
the double-condition subset-product problem, and the density gap itself as a
diagnostic applicable to other Carmichael-type constructions under a double
side condition.
\end{enumerate}

\emph{What is not claimed.} This paper does \textbf{not} disprove Agrawal's
conjecture, which remains open (and widely believed false). It does
\textbf{not} prove that no counterexample can ever be constructed, nor even
that the L--P construction is infeasible in principle. It establishes,
quantitatively, that one specific and natural pipeline --- twin-smooth
harvest followed by subset-product solve at the AKS modulus $r = 5$, with
the known solver families --- is blocked at the solver stage, and it locates
the blockage precisely: a measured identity-density threshold
($\rho \gtrsim 0.98$) that no harvestable pool under the double condition
comes within two orders of magnitude of meeting. What remains open is
exactly what the measurements do not touch: solver paradigms outside the
linear-algebra/birthday families, and constructions at a different modulus
$r$; see \S\ref{sec:scope}, which should be read as part of the result, not
as a disclaimer.

\emph{Provenance of numbers.} Every quantitative claim below is a
\emph{measurement} produced by a script in the companion repository, recorded
in its findings log; Appendix~\ref{app:repro} maps each claim to the script
and data file that produced it. Throughout we distinguish \emph{proved}
statements (the lemmas of \S\ref{sec:construction}), \emph{measured}
quantities (ratios, densities, solver frontiers), and \emph{heuristic or
conjectural} inputs (the L--P counterexample-density heuristic, Dickman-type
yield extrapolations). The ``blocked'' conclusion is an empirical and
structural argument, not a theorem.

\section{The construction and its reduction}\label{sec:construction}

The statements in this section are proved in the companion working paper
\cite{WP}; we restate them without proof, except where the proof is one line.

\begin{theorem}[Lenstra--Pomerance \cite{LPremarks}; statement following
V\'a\v{n}a \cite{Vana09}, who proves the $k \equiv 3 \pmod 4$ case]
\label{thm:lp}
Let $n = p_1 \cdots p_k$ with $p_1, \dots, p_k$ distinct primes such that
\begin{enumerate}\itemsep1pt
\item[(a)] $k$ is odd;
\item[(b)] $p_i \equiv 3 \pmod{80}$ for all $i$;
\item[(c)] $p_i - 1 \mid n - 1$ for all $i$ \hfill ($n$ is Carmichael);
\item[(d)] $p_i + 1 \mid n + 1$ for all $i$ \hfill ($n$ is Lucas--Carmichael).
\end{enumerate}
Then $T(-1,n,5)$ holds and $n^2 \not\equiv 1 \pmod 5$; in particular such an
$n$ with $k \ge 3$ is a counterexample to Conjecture~\ref{conj:agrawal}.
\end{theorem}

The entire difficulty of satisfying (c) and (d) simultaneously is
concentrated in the following trivial but decisive observation.

\begin{lemma}[Side-coprimality]\label{lem:sides}
If $n$ satisfies (c) and (d), then $\gcd(p_i-1,\,p_j+1) \mid 2$ for all $i,j$
(including $i=j$). Consequently the odd primes dividing
$\prod_i (p_i-1)$ and those dividing $\prod_i (p_i+1)$ form disjoint sets.
\end{lemma}
\begin{proof}
An odd prime $q$ dividing $p_i-1$ and $p_j+1$ divides both $n-1$ and $n+1$,
hence their difference $2$, a contradiction.
\end{proof}

Under condition (b) the $2$-adic and mod-$5$ parts of (c),(d) are automatic
\cite[Lem.~2.5]{WP}, so the problem lives entirely in odd moduli prime to
$5$, split by Lemma~\ref{lem:sides} into two disjoint prime universes --- a
``minus side'' supporting the $p_i - 1$ and a ``plus side'' supporting the
$p_i + 1$.

\subsection{Reduction to a two-sided subset-product problem}

Fix disjoint finite sets $Q_-, Q_+$ of odd primes with
$5 \notin Q_- \cup Q_+$, and define the \emph{pool}
\[
\cP \;=\; \Bigl\{\, p \ \text{prime} : p \equiv 3 \!\!\pmod{80},\;
\tfrac{p-1}{2} \text{ is } Q_-\text{-smooth},\;
\tfrac{p+1}{4} \text{ is } Q_+\text{-smooth} \Bigr\}.
\]
Let $\Lambda_- = 2\,\lcm\{p-1 : p \in \cP\}$ and
$\Lambda_+ = 4\,\lcm\{p+1 : p \in \cP\}$.

\begin{proposition}[{\cite[Prop.~3.1]{WP}}]\label{prop:reduction}
Let $S \subseteq \cP$ with $|S|$ odd, $|S| \ge 3$, and $n = \prod_{p \in S} p$.
If
\begin{equation}\label{eq:target}
n \equiv 1 \pmod{\Lambda_-} \qquad\text{and}\qquad n \equiv -1 \pmod{\Lambda_+},
\end{equation}
then $n$ satisfies conditions (a)--(d) of Theorem~\ref{thm:lp}, hence is a
counterexample to Conjecture~\ref{conj:agrawal}.
\end{proposition}

Condition \eqref{eq:target} is stronger than (c),(d) --- it imposes the full
lcm rather than the lcm over $S$ --- but in exchange it is linear: taking
discrete logarithms in each cyclic component of
$G = (\Z/\Lambda_-)^\times \times (\Z/\Lambda_+)^\times$ (Pohlig--Hellman
\cite{PH78} is cheap, all component orders being smooth by design),
\eqref{eq:target} becomes a $0/1$ subset-sum over $\bigoplus_i \Z/m_i$, with
one extra $\Z/2$ coordinate enforcing $|S|$ odd. The heuristic solvability
condition is $2^{|\cP|} \gg |G|$, i.e.\
$|\cP| \gtrsim \log_2 \Lambda_- + \log_2 \Lambda_+$ (the \emph{demand bits})
with slack for the solver.

\subsection{The twin-smooth dictionary}

\begin{lemma}[Dictionary; {\cite[Lem.~4.1]{WP}}]\label{lem:dict}
Let $m \ge 1$ and $p = 2m+1$. Then $p \in \cP$ (for $Q_\pm$ large enough to
contain the relevant primes) if and only if
\begin{enumerate}\itemsep1pt
\item[(i)] $m$ and $m+1$ are both smooth with respect to
$Q_- \cup Q_+ \cup \{2\}$, with the odd part of $m$ being $Q_-$-smooth and
the odd part of $m+1$ being $Q_+$-smooth;
\item[(ii)] $m \equiv 1 \pmod{40}$;
\item[(iii)] $p$ is prime.
\end{enumerate}
Moreover (ii) alone implies $p \equiv 3 \pmod{80}$, $v_2(p-1)=1$,
$v_2(p+1)=2$, and $5 \nmid m(m+1)$.
\end{lemma}

Pairs of consecutive smooth integers $(m, m+1)$ --- \emph{twin smooths} ---
with $2m+1$ prime are precisely the parameters sought for the isogeny-based
schemes B-SIDH \cite{BSIDH} and SQIsign, and 2020--2025 produced an
industrial literature on generating them: PTE sieving \cite{CMN21,Sterner23},
the optimized Conrey--Holmstr\"om--McLaughlin (CHM) iteration
\cite{CHM13,BCCEMNS23} (whose $B=547$ run yields $82{,}026{,}426$ pairs),
Pell/St\o rmer complete sets \cite{Stormer,Lehmer64,BHLNV}, and
lattice-based record hunts \cite{MSW}. By Lemma~\ref{lem:dict}, a congruence
slice of any twin-smooth dataset, filtered for primality of $2m+1$, is a
ready-made pool --- before the one condition that literature never needs: the
\emph{global} side-partition forced by Lemma~\ref{lem:sides}.

\subsection{Pipeline vocabulary}\label{sec:pipeline}

The repository catalogs the pipeline components in a registry
(\texttt{doc/registry.md}) whose identifiers we adopt: harvesters
\textbf{H1}--\textbf{H8} (H1 native sieve; H2/H3 CHM closure in Python/C; H4
PTE import; H5 corpus streaming; \textbf{H6} partition-first; \textbf{H7}
co-designed partition-first; \textbf{H8} divisor-paradigm), solvers
\textbf{S1}--\textbf{S7} (S1 partition/coloring optimizer; \textbf{S2}
GF(2) $2$-part elimination; S3 annealing; S4 exact meet-in-the-middle;
S5 per-prime reduction scaffold; \textbf{S6} AGHS $\omega$-guided reduction;
\textbf{S7} Wagner-$4$ birthday, used as a density-reach instrument),
and validators V1--V4. The narrative of
\S\S\ref{sec:supply}--\ref{sec:solver} is: mining (H1--H5\,$\to$\,S1) fails
by a measured invariant; partition-first harvest (H6/H7) solves supply;
solvers S2--S4 hit a dimensional frontier; the scalable S6 needs a pool
structure (H8) that the double condition makes unharvestable; and the
reach measurement S7 shows the low-density pools that \emph{are}
harvestable lie beyond the entire birthday family's reach.

\section{Supply is solved}\label{sec:supply}

This section establishes the positive half: pool supply is not the binding
constraint. All quantities are measurements; scripts and data files are
listed in Appendix~\ref{app:repro}.

\subsection{Mining fixed corpora is dead: the $0.06$ invariant}
\label{sec:invariant}

The obvious first move is to mine the existing twin-smooth corpora: take all
twins $(m,m+1)$ below a smoothness bound $B$, keep the dictionary slice
($m \equiv 1 \pmod{40}$, $2m+1$ prime), and then choose the side-partition
$Q_- \sqcup Q_+$ maximizing the number of \emph{surviving} elements ---
those whose $m$-support lies entirely in $Q_-$ and $(m+1)/2$-support in
$Q_+$. Feasibility requires survivors comparable to the demand bits of the
group they generate.

\begin{finding}[The $0.06$ mining invariant --- measured]\label{find:invariant}
Across seven experiments spanning $B = 100$--$547$, truncated and complete
corpora, and pool sizes $78$ to $73{,}163$, the optimized partition retains
\[
\frac{\#\text{survivors}}{\#\text{demand bits}} \;\approx\; 0.060\text{--}0.083,
\]
with element-greedy optimization and direct partition-space local search
agreeing (Table~\ref{tab:invariant}). Feasibility needs this ratio
$\gtrsim 1$; the shortfall is a factor $16$--$25$.
\end{finding}

The measurement includes the \emph{complete} $B=547$ corpus of
$82{,}026{,}426$ pairs (shared by M.~Naehrig; $642{,}495$ pairs with
$m \equiv 1 \pmod{40}$, of which $73{,}163$ have $2m+1$ prime): the best
partition found keeps $82$ survivors against $984$ demand bits (ratio
$0.083$). The invariant is thus measured, not extrapolated, on the entire
published supply. The mechanism is visible in the data: partition-agnostic
twins carry a mean total of $\approx 15$ distinct odd primes across both
sides at $B = 547$, so each kept element ``spends'' roughly two fresh primes
($\approx 15$ bits of new demand) that no later element may straddle.
Optimization beats a random partition by $3.4$--$4.4\times$, but no
optimizer can beat the arithmetic: \textbf{dataset mining is dead as a
primary route}.

\subsection{Partition-first harvesting (H6) and its gates}
\label{sec:partitionfirst}

The invariant is a property of \emph{partition-agnostic} supply. Inverting
the order --- fix the partition \emph{first}, then harvest only compatible
elements --- removes it by construction. Concretely (harvester H6): choose
$Q_+$ = the odd primes $\le B$ (except $5$) and $Q_-$ = a band of medium
primes in $(B, B']$; enumerate $m$ as products of $j$ primes of $Q_-$
restricted to $m \equiv 1 \pmod{40}$; per candidate, trial-divide $(m+1)/2$
over $Q_+$ and test $p = 2m+1$ for primality. Every hit is compatible with
the partition by construction; the cost moves from a global combinatorial
obstruction into per-candidate yield, a local Dickman-type probability.

Two pre-registered gates governed the program. \textbf{G1} (yield exists):
some measured cell must show pool/demand $\ge 0.5$ at prototype scale with a
fitted extrapolation $\ge 2.5$. \textbf{G2} (production spec): a frozen
specification with projected pool/demand $\ge 2.5$ and $\ge 3000$ usable
elements. Both passed:

\begin{finding}[Partition-first supply --- measured]\label{find:supply}
On the measured $j=4$ grid (Table~\ref{tab:ratios}), pool-to-demand ratios
reach $13.8$ at $(B, B') = (1259, 12590)$ --- the frozen production
specification: $Q_+ =$ odd primes $\le 1259$ ($\neq 5$), $Q_- \subset
(1259, 12590]$, $j = 4$, projected pool $\approx 2.5\times 10^5$ elements of
mean $50$ bits against $\approx 18{,}100$ demand bits --- and up to
$41$--$45.5$ at larger $B'$. Partition-first thus beats the mining invariant
by a factor of roughly $200$--$700$ in the best cells.
\end{finding}

\begin{finding}[Model validation --- measured, conservative]
\label{find:validation}
The sampling-based yield model behind Table~\ref{tab:ratios} was validated
out-of-sample against the complete $B=547$ corpus (exhaustive twin
population to $122$ bits): across four asymmetric splits it predicted
$0.41$--$1.04\times$ the true pool counts (geometric mean $0.68$) --- i.e.\
the model is conservative, never optimistic, and the contrast with the
mining route ($\approx 300$--$500\times$ separation) is measured at scale on
both sides with the same congruence, filters, and optimizer.
\end{finding}

\subsection{Corpus and closure validation}\label{sec:closure}

Supporting the supply claims, the repository's C implementation of the CHM
closure (H3, \texttt{chm\_closure.c}) was validated against ground truth:
it recovers $13{,}333$ of the complete $13{,}374$ pairs at $B = 100$
($99.7\%$), $33{,}118$ at $B=113$ ($99.65\%$ of the complete $33{,}233$),
and $346{,}110$ at $B = 200$ --- $99.98\%$ of the originally published
$346{,}192$ and $99.2\%$ of the true complete count $348{,}840$ (the
published figure was itself $0.8\%$ short), with zero false positives and
exact agreement on the largest pair in each case. The anchored counts
$N(100) = 13{,}374$, $N(200) = 348{,}840$, $N(547) = 82{,}026{,}426$ fit
$N(B) \sim B^{\alpha}$ with $\alpha \approx 4.7$--$5.4$, so fixed-$B$
supply, while finite (St\o rmer \cite{Stormer,Lehmer64}), grows steeply in
$B$.

\medskip
\noindent\textbf{Conclusion of \S\ref{sec:supply}.} Supply is abundant,
validated, and over-provisioned by one to two orders of magnitude relative
to the heuristic feasibility threshold. Whatever blocks the pipeline is
downstream.

\section{The solver obstruction}\label{sec:solver}

By Proposition~\ref{prop:reduction} the remaining task is: given the pool,
find $S$ with $|S|$ odd and $\prod_{p\in S} p \equiv (1, -1)$ in
$G = (\Z/\Lambda_-)^\times \times (\Z/\Lambda_+)^\times$. Via CRT, $G$
decomposes into cyclic prime-power components $\Z/\ell^a$ arising from the
prime-power parts of each $q - 1$, $q \in Q_- \cup Q_+$; Pohlig--Hellman
discrete logarithms turn the target into a $0/1$ subset-sum, NP-hard in
general and tractable only through structure. With pool size
$N \gg \dim G$, solutions are heuristically super-dense; \emph{finding} a
$0/1$ one is the crux.

\subsection{The $2$-part is easy --- and small}\label{sec:twopart}

Every component reduced mod $2$ is linear over $\mathrm{GF}(2)$ (the
unknowns being $0/1$), so the $2$-part of the system, plus the parity row,
solves by Gaussian elimination. This scales without difficulty:

\begin{finding}[Measured]\label{find:gf2}
At the frozen specification $(B, B') = (1259, 12590)$ the group has $4586$
cyclic components and the full $\mathrm{GF}(2)$ system solves in $12$
seconds (solver S2). But the $2$-part accounts for only $\approx 16\%$ of
the demand bits.
\end{finding}

This overturned the program's working assumption that the $2$-torsion was
the bulk of the problem. The remaining $\approx 84\%$ ---
$\approx 15{,}000$ bits at the frozen spec --- lives in odd-prime
components with primes $\ell$ up to $6269$.

\subsection{The odd part defeats the generic solvers}\label{sec:oddpart}

Two generic attacks on the odd part were built and pushed to their limits on
faithful synthetic instances with planted solutions
(Table~\ref{tab:frontier}): simulated annealing inside the
$\mathrm{GF}(2)$ solution family (S3) stalls at $\ge 19$ odd components,
and exact meet-in-the-middle (S4) solves what annealing cannot
($N{=}40$, $15$ odd components, in $11$s) but is exponential in $N$, capping
near $N \approx 40$ pool elements, $\approx 15$ odd components,
$\approx 40$ odd bits. (Wagner's generalized-birthday algorithm
\cite{Wagner}, the other natural generic candidate, needs list sizes
exponential in the odd bit-count in its classic form and was not
competitive in this regime.)

A harvest/solver co-design (H7) attacks the odd part from the supply side:
constraining both factor bases to primes $q$ whose $q-1$ has odd part
$t_0$-smooth caps the largest odd prime in the group. Measured co-designed
cells (Table~\ref{tab:codesign}) show a viable window --- at
$t_0 = 100$, $B = 2003$, $B' = 60090$ the harvest ratio stays $7.7$ (above
the G2 gate) while the odd support shrinks from $336$ primes $\le 6269$ to
$24$ primes $\le 97$. But co-design caps the odd prime \emph{size}, not the
odd \emph{dimension}: recovering yield requires many primes, and the odd
part remains at $\approx 24{,}000$ bits across thousands of components ---
two to three orders of magnitude beyond the S3/S4 frontier. Squeezing the
dimension itself would require near-Fermat ($2$-power-heavy) primes, which
are too sparse to harvest: a genuine yield/solver tension.

\begin{finding}[Measured]\label{find:frontier}
Generic exact and heuristic solvers for the odd part cap at $\approx 15$ odd
components / $\approx 40$ odd bits, versus $\approx 24{,}000$ odd bits
(thousands of components) required at the co-designed specification --- a
shortfall of two to three orders of magnitude that co-design provably (by
measurement) does not close.
\end{finding}

\subsection{The scalable method and what it requires}\label{sec:aghs}

The frontier of Finding~\ref{find:frontier} is not the end of the story,
because subset-products $\equiv 1$ in $(\Z/\Lambda)^\times$ at vastly larger
scales \emph{have} been solved: Alford, Grantham, Hayman, and Shallue
\cite{AGHS}, building on L\"oh--Niebuhr \cite{LN} and the AGP machinery
\cite{AGP}, constructed a Carmichael number with $10{,}333{,}229{,}505$
prime factors. Their method is not linear algebra: it is a subexponential
$\omega$-guided reduction, driving a distance-to-identity potential (the
highest CRT component at which a partial product differs from the identity)
to zero top-down, with a birthday step. It has no small-odd-part
requirement, which is precisely why the generic-solver frontier above is a
statement about the wrong solver family rather than about the problem.

The method does, however, have a structural input requirement. AGHS pools
consist of primes $p$ with $p - 1 \mid \Lambda$ for a fixed highly composite
$\Lambda$: such a $p$ is $\equiv 1$ on \emph{every} component whose prime
power divides $p-1$, i.e.\ the pool vectors are identity on a large fraction
of the components. Call this fraction the \emph{identity density} of a pool
element. High identity density is what makes $\omega$-reduction move: most
steps must leave most components untouched. AGHS-style pools have identity
density $\approx 0.3$--$0.5$. Our partition-first pools are the opposite
extreme: $p - 1 = 2m$ with $m$ a product of $j \approx 4$ medium primes is
identity on only $\approx j$ of the components touched by the pool ---
measured density $\approx 0.003$--$0.008$, about $100\times$ lower.

\subsection{The harvest jaw: the divisor paradigm is unharvestable}
\label{sec:gap}

If the double condition admitted an AGHS-style harvest, the pipeline would
be back on track. The natural design (harvester H8, the \emph{divisor
paradigm}) fixes disjoint factor bases $T_-, T_+$ (side-coprimality), sets
$\Lambda_-$, $\Lambda_+$ highly composite over them, and harvests pool
primes $p = 2m+1$ with $m$ a (squarefree) divisor of $\prod T_-$,
$m \equiv 1 \pmod{40}$, $2m+1$ prime, and $(m+1)/2$ squarefree over $T_+$
--- exactly the $p - 1 \mid \Lambda_-$, $p+1 \mid \Lambda_+$ structure the
$\omega$-solver wants, on both sides at once.

\begin{finding}[N1: the divisor paradigm is unharvestable --- measured]
\label{find:n1}
Across all feasible scales tested (Table~\ref{tab:n1}), the H8 pool
contains $0$--$2$ primes. Small $T_-$ (needed for decent density
$\omega(m)/|T_-|$) yields an empty or near-empty pool; enlarging $T_-$
collapses the density ($0.178 \to 0.067$--$0.088$) while the pool stays at
$1$--$2$. The root cause is structural: the double condition ---
\emph{both} $p - 1 \mid \Lambda_-$ \emph{and} $p + 1 \mid \Lambda_+$ for
fixed moduli --- is the product of two rare events, far more restrictive
than the single condition $p - 1 \mid \Lambda$ that AGHS harvests satisfy.
\end{finding}

\subsection{The solver jaw: measured birthday reach at low density}
\label{sec:reach}

The finding above leaves one live question: perhaps the scalable solver
family does not actually \emph{need} AGHS-like density, and would tolerate
the $\rho \approx 0.003$--$0.008$ of the harvestable H7 pools. This was
tested directly. A correct Wagner-$4$ generalized-birthday solver
\cite{Wagner} --- a four-list $4$-sum to the identity over a small-prime
component group, the collision engine at the core of the AGHS-class methods
--- was implemented (S7) and its \emph{reach} measured: the largest
component count $r_{\max}$ it solves, as a function of identity density
$\rho$ and pool size $N$, on synthetic pools whose coordinates are identity
with probability $\rho$.

\begin{finding}[A10: the reach law --- measured]\label{find:reach}
The measured reach obeys
\[
(1-\rho)\, r_{\max} \;\approx\; \mathrm{const}
\qquad\text{with}\qquad r_{\max} \;\sim\; \log N
\]
(Table~\ref{tab:reach}): at H7's density $\rho \approx 0.004$ the solver
reaches $r_{\max} = 20$ components at $N = 20{,}000$ and only $24$ at
$N = 100{,}000$, while at $\rho = 0.9$ it reaches $128$. Our real groups
have \emph{thousands} of components. Two consequences:
\begin{enumerate}\itemsep1pt
\item[(i)] \textbf{Density threshold.} Reaching a $\sim$$1000$-component
group needs $(1-\rho)\cdot 1000 \lesssim 20$, i.e.\ $\rho \gtrsim 0.98$
(optimistic higher-$k$ Wagner variants lower this only to
$\approx 0.85$--$0.95$). H7 supplies $\rho \approx 0.004$ --- two-plus
orders of magnitude short.
\item[(ii)] \textbf{Pool futility.} Since $r_{\max} \sim \log N$,
harvesting more cannot close the gap: a $5\times$ larger pool bought $+4$
components, and the group itself grows with the harvest, so the ratio
never improves.
\end{enumerate}
\end{finding}

\subsection{The two-part obstruction (the core negative result)}
\label{sec:vise}

\begin{finding}[The identity-density vise --- measured]\label{find:gap}
Under the double Carmichael + Lucas--Carmichael condition:
\begin{center}\footnotesize
\begin{tabular}{@{}llll@{}}
\toprule
pool & identity density & pool size & birthday-solvable? \\
\midrule
AGHS-style ($p-1 \mid \Lambda$, single condition) & $\approx 0.3$--$0.5$ &
large & yes (their record) \\
partition-first H7 (harvestable) & $\approx 0.003$--$0.008$ &
$\approx 10^5$ & no (needs $\rho \gtrsim 0.98$) \\
divisor-paradigm H8 (AGHS density) & $\approx 0.07$--$0.18$ if nonempty &
$0$--$2$ & moot \\
\bottomrule
\end{tabular}
\end{center}
One can have a harvestable pool \emph{or} solver-reachable identity density
--- not both: high-density pools are unharvestable
(Finding~\ref{find:n1}), and low-density pools are unsolvable
(Finding~\ref{find:reach}), with no other family in the measured frontier
(Table~\ref{tab:frontier}) reaching the required dimension. That is the
obstruction.
\end{finding}

This is the punchline of the paper. The supply half of the pipeline is
solved and over-provisioned (\S\ref{sec:supply}); the solve half is closed
for the entire AGHS/birthday solver class, not by an absence of solvers,
but by a two-part quantified mismatch: the class needs identity density
$\rho \gtrsim 0.98$, the double condition harvests at
$\rho \approx 0.004$, and the two measurements bounding the gap ---
harvest yield of high-density pools, and solver reach on low-density
pools --- close it from both sides. In particular the $r = 5$ (mod-$80$)
Lenstra--Pomerance route, as a harvest-then-solve pipeline, fails not for
lack of supply but because its double side condition structurally caps the
identity density its solvers would need.

\section{Scope, caveats, and what remains open}\label{sec:scope}

This section bounds the result. It is part of the result.

\begin{enumerate}\itemsep3pt
\item \textbf{Not a disproof of Agrawal's conjecture.} The conjecture
remains open and is believed false; nothing here bears on its truth. A
reviewer's first question --- ``did you prove the conjecture has no
counterexample?'' --- has answer \emph{no}, emphatically.
\item \textbf{Not a proof that no counterexample is constructible.} The
blockage established is of \emph{this pipeline} (twin-smooth harvest
$\to$ subset-product solve at $r = 5$) with the \emph{known solver
families} (linear-algebraic, annealing, meet-in-the-middle, and the
AGHS/birthday class as characterized by its measured density-reach law).
Other constructions, and solvers outside these families, are unconstrained
by our measurements.
\item \textbf{Finite-scale and squarefree caveats on N1.}
Finding~\ref{find:n1} tested squarefree divisors at finite scales
(Table~\ref{tab:n1}). Prime-power moduli $\Lambda_\pm$ are untested,
though the structural analysis points the same way (high density requires
$m$ to retain full prime powers, forcing $m$ large); enormous scales are
untested, and the measured density trend is against them.
\item \textbf{Caveats on the reach law (A10).}
Finding~\ref{find:reach} measures Wagner-$4$, the simplest birthday
variant, on a model group with uniform $\Z/3$ components; real components
are mixed small primes, for which the same qualitative law is expected but
was not separately measured. Higher-$k$ Wagner variants reach more
components ($\approx$ linearly in the level count $h$, but $h \le \log_2 N$
and each level costs pool), which lowers the density threshold only to
$\rho \approx 0.85$--$0.95$ --- still roughly $100\times$ above what H7
supplies. The complete AGHS algorithm (two potentials, non-uniformity
exploitation) was not reproduced in full; the closure of the low-density
escape rests on the measured reach law of the birthday core that the whole
class shares, plus this higher-$k$ analysis. (An earlier partial
$\omega$-reducer, \texttt{a8\_omega\_solver.py}, solves only $r \le 20$
components at high identity density, consistent with the law.)
\item \textbf{What remains open.} Exactly two directions are untouched by
these measurements. \emph{(b) A fundamentally different solver paradigm:}
something outside the linear-algebra, local-search, meet-in-the-middle, and
birthday/$\omega$-reduction families, purpose-built for low-density
many-prime pools. \emph{(c) A different construction --- in particular a
different AKS modulus $r$:} the entire quantitative case here is specific
to the $r = 5$ (mod-$80$) cell --- its congruence slice, its group
$(\Z/\Lambda_-)^\times \times (\Z/\Lambda_+)^\times$, and its double side
condition; a different $r$ changes all three, and some $r \neq 5$ may admit
denser pools or a friendlier group. Related is the admissible-cell question
of \cite[Q.~6.2]{WP}: the boost obstruction \cite[Prop.~6.1]{WP} shows the
mod-$80$ cell is unreachable by the twin-smooth boosting family $2x^n - 1$
\cite{BCCEMNS23,Sterner23}, but whether \emph{other} local-condition cells
admit the L--P mechanism (and are reachable by structurally side-separated
families) is open, and a positive answer would restructure the harvest
entirely.
\end{enumerate}

\section{Reusable contributions and conclusion}\label{sec:conclusion}

Beyond the negative result, the program leaves behind machinery of
independent use:

\begin{itemize}\itemsep2pt
\item \textbf{A validated twin-smooth supply model.} The partition-first
yield model, validated conservative (geometric mean $0.68\times$,
never optimistic) against the complete $B = 547$ corpus, plus the
near-complete C implementation of the CHM closure ($99.2$--$99.98\%$ of the
exact complete sets, zero false positives) --- including the observation
that the published $N(200) = 346{,}192$ was itself $0.8\%$ short of the
true $348{,}840$.
\item \textbf{The $0.06$ mining invariant} (Finding~\ref{find:invariant}):
a sharp, reproducible measurement of the global side-partition cost for
partition-agnostic twin corpora, and its complete bypass by
partition-first harvesting --- relevant to any application that must
multiplicatively combine twins under a side constraint.
\item \textbf{A solver-frontier map and reach law} for two-sided $0/1$
subset-product problems (Tables~\ref{tab:frontier} and~\ref{tab:reach}),
with a synthetic-instance regression harness (planted solutions, tunable
identity density $\rho$) against which any proposed new solver paradigm can
be tested down to the harvestable regime $\rho \approx 0.004$ --- and the
reach law $(1-\rho)\,r_{\max} \approx \mathrm{const}$, $r_{\max} \sim \log
N$, as the baseline any such paradigm must beat.
\item \textbf{The identity-density vise as a diagnostic.} For any
Carmichael-type constructive problem under a \emph{double} side condition,
three measurable quantities --- harvest yield of divisor-structured pools,
identity density of harvestable pools, and birthday reach at that density
--- locate the pipeline on the harvestable/solvable dichotomy of
Finding~\ref{find:gap} before major compute is spent.
\end{itemize}

\emph{Conclusion.} The $r = 5$ Lenstra--Pomerance construction reduces
cleanly to a two-sided subset-product problem whose supply side is,
contrary to initial expectation, solved: partition-first harvesting
over-provisions the pool by one to two orders of magnitude. The pipeline
is closed at the solver stage by a two-part quantified obstruction. The
scalable AGHS/birthday solver class obeys a measured reach law
$(1-\rho)\,r_{\max} \approx \mathrm{const}$ with $r_{\max} \sim \log N$,
so our thousand-component groups demand identity density
$\rho \gtrsim 0.98$; the double Carmichael + Lucas--Carmichael condition
caps harvestable pools at $\rho \approx 0.003$--$0.008$, and the
divisor-paradigm harvest that would supply solver-grade density is empty
at every feasible scale. High-density pools are unharvestable; low-density
pools are unsolvable; larger harvests help only logarithmically. Within
the known solver families, the constructive route to refuting Agrawal's
conjecture via twin smooths at $r = 5$ is therefore a measured dead end:
\emph{supply yes, solver no, and the wall is quantified from both sides}.
What survives is precisely what the measurements do not reach --- a
genuinely different solver paradigm, or a different modulus $r$.

\appendix

\section{Data tables}\label{app:tables}

All tables reproduce measurements from the repository findings log
\texttt{data/FINDINGS.md}; the generating script and log section are noted
under each table.

\begin{table}[ht]\centering\small
\begin{tabular}{@{}llrrl@{}}
\toprule
$B$ & corpus & pool & optimized survivors & surv/demand-bits \\
\midrule
113 & sieve, $X \le 10^8$ & 59 & 8 & 0.061 \\
113 & sieve, $X \le 10^{10}$ & 73 & 8 & 0.061 \\
113 & CHM closure (complete-ish) & 78 & 8 & 0.061 \\
257 & sieve, $X \le 10^9$ & 497 & 21 & 0.061 \\
547 & sieve, $X \le 10^9$ & 2{,}715 & 51 & 0.068 \\
200 & CHM closure (near-complete) & 562 & 17 / 16 & 0.060 / 0.061 \\
547 & complete corpus (82{,}026{,}426 pairs) & 73{,}163 & 82 & 0.083 \\
\bottomrule
\end{tabular}
\caption{The $0.06$ mining invariant (Finding~\ref{find:invariant}).
``Pool'' = twins with $m \equiv 1 \pmod{40}$ and $2m+1$ prime; survivors
under the best side-partition found (element-greedy; for $B=200$ also
direct partition-space local search, second figure). Scripts:
\texttt{phase1a\_calibration.py}, \texttt{chm\_closure.c},
\texttt{b547\_pool\_analysis.py}; FINDINGS ``Calibration runs'',
``Compiled CHM port'', ``Definitive $B=547$ coloring measurement''.}
\label{tab:invariant}
\end{table}

\begin{table}[ht]\centering\small
\begin{tabular}{@{}rrrrrr@{}}
\toprule
$B$ & $B'$ & ratio & pool & demand bits & rel.\ err \\
\midrule
1259 & 12590 & 13.8 & $2.5\times 10^5$ & 18{,}117 & $\pm 24\%$ \\
1259 & 18885 & 41.0 & $1.1\times 10^6$ & 27{,}101 & $\pm 26\%$ \\
2003 & 20030 & 45.5 & $1.3\times 10^6$ & 28{,}768 & $\pm 24\%$ \\
2003 & 40060 & 209\phantom{.0} & $1.2\times 10^7$ & 57{,}514 & $\pm 32\%$ \\
\bottomrule
\end{tabular}
\caption{Partition-first (H6) measured pool/demand ratios, reliable $j=4$
cells (Finding~\ref{find:supply}); the $(1259, 12590)$ row is the frozen
specification. Script: \texttt{a2\_grid.py} (data
\texttt{a2\_grid.json}); FINDINGS ``A2: model validated \dots\ G2
passed''.}
\label{tab:ratios}
\end{table}

\begin{table}[ht]\centering\small
\begin{tabular}{@{}llll@{}}
\toprule
solver & method & reaches & caps at \\
\midrule
S2 & GF(2) elimination ($2$-part + parity) & 4586 components, 12 s &
$2$-part only ($\approx 16\%$ of bits) \\
S3 & annealing in GF(2) solution family & --- & stalls $\ge 19$ odd
components \\
S4 & exact meet-in-the-middle & $N{=}40$, $15$ odd comps, 11 s &
$\approx 40$ odd bits (exp.\ in $N$) \\
\midrule
\multicolumn{2}{@{}l}{required (co-designed spec)} &
\multicolumn{2}{l@{}}{$\approx 24{,}000$ odd bits, thousands of odd
components} \\
\bottomrule
\end{tabular}
\caption{The generic solver frontier on the odd part
(Finding~\ref{find:frontier}), on faithful synthetic instances with planted
solutions. Scripts: \texttt{a4\_solver.py}, \texttt{a6\_oddsolver.py};
FINDINGS ``A4 solver, G3 phase-1'' and ``A6 small-prime odd-part
solver''.}
\label{tab:frontier}
\end{table}

\begin{table}[ht]\centering\small
\begin{tabular}{@{}rrrrrr@{}}
\toprule
$t_0$ & $B$ & $B'$ & ratio & max odd $\ell$ & \# distinct odd \\
\midrule
50 & 1259 & 40060 & 1.1 & 47 & 14 \\
100 & 1259 & 40060 & 2.6 & 97 & 24 \\
\textbf{100} & \textbf{2003} & \textbf{60090} & \textbf{7.7} & \textbf{97}
& \textbf{24} \\
200 & 1259 & 40060 & 14.5 & 199 & 45 \\
50 & 631 & 40060 & 0.0 & 47 & 14 \\
\bottomrule
\end{tabular}
\caption{Co-designed harvest (H7): measured yield with both factor bases
constrained to odd part of $q - 1$ being $t_0$-smooth (hit counts low,
$\pm 40\%$). Bold row: the viable window. Scripts:
\texttt{a5\_codesign.py}, \texttt{a5\_yield.py} (data
\texttt{a5\_codesign\_yield.json}); FINDINGS ``A5 harvest/solver
co-design''.}
\label{tab:codesign}
\end{table}

\begin{table}[ht]\centering\small
\begin{tabular}{@{}rrrrrr@{}}
\toprule
$T_-$ bound & $T_+$ range & $|T_-| + |T_+|$ & pool & mean density &
max density \\
\midrule
90 & 97--220 & 45 & 1 & 0.178 & 0.178 \\
120 & 127--300 & 60 & 1 & 0.067 & 0.067 \\
150 & 160--360 & 68 & 2 & 0.074 & 0.088 \\
60 & 67--160 & 34 & 0 & --- & --- \\
\bottomrule
\end{tabular}
\caption{N1: divisor-paradigm harvest (H8) yield and identity density
under the double condition (Finding~\ref{find:n1}). Script:
\texttt{a9\_divisor\_harvest.py}; FINDINGS ``N1: divisor-paradigm harvest
(H8) is unharvestable''.}
\label{tab:n1}
\end{table}

\begin{table}[ht]\centering\small
\begin{tabular}{@{}lrrr@{}}
\toprule
$\rho$ & $N$ & $r_{\max}$ & $(1-\rho)\, r_{\max}$ \\
\midrule
\textbf{0.004 (H7)} & 20{,}000 & \textbf{20} & 19.9 \\
\textbf{0.004 (H7)} & 100{,}000 & \textbf{24} & 23.9 \\
0.5 & 20{,}000 & 28 & 14.0 \\
0.9 & 20{,}000 & 128 & 12.8 \\
\bottomrule
\end{tabular}
\caption{A10: measured Wagner-$4$ birthday reach --- the largest solvable
component count $r_{\max}$ as a function of identity density $\rho$ and
pool size $N$ (Finding~\ref{find:reach}). The product $(1-\rho)\,r_{\max}$
is roughly constant and $r_{\max}$ grows only $\sim \log N$; real groups
have thousands of components. Script: \texttt{a10\_wagner\_density.py};
FINDINGS ``A10 (option a): the low-density escape is closed''.}
\label{tab:reach}
\end{table}

\begin{table}[ht]\centering\footnotesize
\begin{tabular}{@{}llll@{}}
\toprule
pool structure & identity density & harvestable? & solvable? \\
\midrule
AGHS-style ($p - 1 \mid \Lambda$) &
$\approx 0.3$--$0.5$ & no (Tab.~\ref{tab:n1}: pool $0$--$2$) & --- \\
partition-first H7 &
$\approx 0.003$--$0.008$ & yes (pool $\approx 10^5$) &
no (needs $\rho \gtrsim 0.98$, Tab.~\ref{tab:reach}) \\
\bottomrule
\end{tabular}
\caption{The identity-density vise (Finding~\ref{find:gap}): the double
condition supplies harvestable pools (H7: $p - 1 = 2m$ with $m$ a product
of $j \approx 4$ medium primes) only $\approx 100\times$ below the density
the AGHS/birthday solver class exploits ($p - 1 \mid \Lambda$, $\Lambda$
highly composite), and the measured reach law shows that class needs
$\rho \gtrsim 0.98$ at our group sizes. FINDINGS ``N1'' (subsection ``The
identity-density gap'') and ``A10''.}
\label{tab:gap}
\end{table}

\section{Reproducibility}\label{app:repro}

The companion repository (\texttt{PRIMES-Conjectures}) contains every
script and committed data file behind the claims of this paper. The
findings log \texttt{data/FINDINGS.md} is the single source of truth for
all quantities; the map below traces each claim to its generator.

\begin{center}\footnotesize
\begin{tabular}{@{}lll@{}}
\toprule
claim & script (\texttt{code/}) & data (\texttt{data/generated/} or log) \\
\midrule
machinery unit tests (T1--T4) & \texttt{verify\_machinery.py} &
FINDINGS ``Phase 0'' \\
$0.06$ invariant, $B \le 257$ & \texttt{phase1a\_calibration.py} &
FINDINGS ``Calibration runs'' \\
CHM closure (C), $B \le 200$ anchors & \texttt{chm\_closure.c} &
\texttt{chm\_B200.log} \\
complete $B=547$ pool + coloring & \texttt{b547\_pool\_analysis.py} &
\texttt{b547\_pool\_stats.json} \\
A1 yield grid (G1) & \texttt{a1\_yield\_prototype.py} &
\texttt{a1\_yield\_grid.json} \\
A2 grid (G2) & \texttt{a2\_grid.py} &
\texttt{a2\_grid.json} \\
A2 out-of-sample validation & \texttt{a2\_validate\_and\_extend.py} &
\texttt{a2\_validation.json} \\
group decomposition, GF(2) solve & \texttt{a4\_solver.py} &
FINDINGS ``A4 solver'' \\
co-design characterization/yield & \texttt{a5\_codesign.py},
\texttt{a5\_yield.py} & \texttt{a5\_codesign\_yield.json} \\
odd-part solver frontier & \texttt{a6\_oddsolver.py} &
FINDINGS ``A6'' \\
$\omega$-reducer (partial S6) & \texttt{a7\_lohniebuhr.py},
\texttt{a8\_omega\_solver.py} & FINDINGS ``A7/A8'' \\
N1 divisor harvest & \texttt{a9\_divisor\_harvest.py} &
FINDINGS ``N1'' \\
A10 Wagner-$4$ density reach (S7) & \texttt{a10\_wagner\_density.py} &
FINDINGS ``A10'' \\
\bottomrule
\end{tabular}
\end{center}

All runs are deterministic (fixed seeds). The complete $B=547$ twin corpus
is private (courtesy of M.~Naehrig) and is \emph{not} in the repository;
only aggregate statistics derived from it are committed
(\texttt{b547\_pool\_stats.json}), which suffices to check every number
quoted here.

\begin{thebibliography}{99}\itemsep1pt\small

\bibitem{AKS} M.~Agrawal, N.~Kayal, N.~Saxena, \emph{PRIMES is in P},
Ann.\ of Math.\ \textbf{160} (2004), 781--793.

\bibitem{BP01} R.~Bhattacharjee, P.~Pandey, \emph{Primality testing}, BTech
thesis, IIT Kanpur, 2001.

\bibitem{LPremarks} H.~W.~Lenstra, Jr., C.~Pomerance, \emph{Remarks on
Agrawal's conjecture}, AIM workshop notes, 2003.
\url{https://aimath.org/WWN/primesinp/articles/html/50a/}

\bibitem{Popovych09} R.~Popovych, \emph{A note on Agrawal conjecture},
Cryptology ePrint 2009/008.

\bibitem{Vana09} T.~V\'a\v{n}a, \emph{Agrawal's conjecture and Carmichael
numbers}, Comenius University Bratislava, 2009.

\bibitem{HD21} A.~Hegde, P.~Devaraj, \emph{Heuristics for the construction of
counterexamples to the Agrawal conjecture}, Springer PROMS \textbf{344}, 2021.

\bibitem{LP2019} H.~W.~Lenstra, Jr., C.~Pomerance, \emph{Primality testing
with Gaussian periods}, J.\ Eur.\ Math.\ Soc.\ \textbf{21} (2019), 1229--1269.

\bibitem{Pinch} R.~G.~E.~Pinch, \emph{The Carmichael numbers up to $10^{18}$},
arXiv:math/0604376.

\bibitem{AGP} W.~R.~Alford, A.~Granville, C.~Pomerance, \emph{There are
infinitely many Carmichael numbers}, Ann.\ of Math.\ \textbf{139} (1994),
703--722.

\bibitem{LN} G.~L\"oh, W.~Niebuhr, \emph{A new algorithm for constructing
large Carmichael numbers}, Math.\ Comp.\ \textbf{65} (1996), 823--836.

\bibitem{AGHS} W.~R.~Alford, J.~Grantham, S.~Hayman, A.~Shallue,
\emph{Constructing Carmichael numbers through improved subset-product
algorithms}, Math.\ Comp.\ \textbf{83} (2014), 899--915; arXiv:1203.6664.

\bibitem{PH78} S.~Pohlig, M.~Hellman, \emph{An improved algorithm for
computing logarithms over $GF(p)$\dots}, IEEE Trans.\ Inf.\ Theory
\textbf{24} (1978), 106--110.

\bibitem{Wagner} D.~Wagner, \emph{A generalized birthday problem}, CRYPTO
2002.

\bibitem{BSIDH} C.~Costello, \emph{B-SIDH: supersingular isogeny
Diffie--Hellman using twisted torsion}, ASIACRYPT 2020.

\bibitem{CMN21} C.~Costello, M.~Meyer, M.~Naehrig, \emph{Sieving for twin
smooth integers with solutions to the Prouhet--Tarry--Escott problem},
EUROCRYPT 2021; ePrint 2020/1283.

\bibitem{BCCEMNS23} G.~Bruno, M.~Corte-Real Santos, C.~Costello,
J.~K.~Eriksen, M.~Meyer, M.~Naehrig, B.~Sterner, \emph{Cryptographic smooth
neighbors}, ASIACRYPT 2023; ePrint 2022/1439.

\bibitem{CHM13} J.~B.~Conrey, M.~A.~Holmstr\"om, T.~L.~McLaughlin,
\emph{Smooth neighbors}, Experiment.\ Math.\ \textbf{22} (2013), 195--202.

\bibitem{BHLNV} J.~Buzek, J.~Hasan, J.~Liu, M.~Naehrig, A.~Vigil,
\emph{Finding twin smooth integers by solving Pell equations},
arXiv:2211.04315.

\bibitem{Sterner23} B.~Sterner, \emph{Towards optimally small smoothness
bounds for cryptographic-sized twin smooth integers and their isogeny-based
applications}, SAC 2024; ePrint 2023/1576.

\bibitem{MSW} E.~Mulder, B.~Sterner, W.~van Woerden, \emph{Large smooth twins
from short lattice vectors}, ANTS XVII; ePrint 2025/1462.

\bibitem{Stormer} C.~St\o rmer, \emph{Quelques th\'eor\`emes sur l'\'equation
de Pell $x^2 - Dy^2 = \pm 1$ et leurs applications}, Skr.\ Norske Vid.\
Akad., 1897.

\bibitem{Lehmer64} D.~H.~Lehmer, \emph{On a problem of St\o rmer}, Illinois
J.\ Math.\ \textbf{8} (1964), 57--79.

\bibitem{WP} D.~Shumow and Claude, \emph{Toward a counterexample to
Agrawal's conjecture via twin smooth integers}, working draft v0.1, August
2026 (companion repository, \texttt{doc/working-paper.tex}).

\end{thebibliography}

\end{document}
